\section{A robust perturbation theorem and the optimized endpoint}
\label{sec:perturbation}

\subsection{A general finite perturbation principle}

The probabilistic step can be separated from the arithmetic geometry.
It applies whenever the exceptional configurations have few
\emph{literal} signatures.

\begin{theorem}[Robust perturbation with bounded key multiplicity]
\label{thm:abstract}
Let \(k,\mu\ge1\) be integers, and let \(\mathcal P\) be a finite
family of ordered \(k\)-tuples of
distinct vertices in a finite set \(X\). Suppose
\(c_0:X\to\Z/s\Z\), \(s\ge2\), and \(\kappa:X\to\cK\) are fixed,
and each key occurs at most \(\mu\) times in any tuple.
Assume each tuple is either:
\begin{enumerate}
 \item rich: for every \(b\in\Z/s\Z\), at least \(\gamma k\)
       entries have outer color different from \(b\); or
 \item exceptional: its literal word
       \(((c_0(x_j),\kappa(x_j)))_{j<k}\) belongs to a family
       of at most \(T\) signatures.
\end{enumerate}
Let \(0<p\le1/2\) and \(t=\fl{pk/2}<\gamma k\). If
\begin{equation}\label{eq:abstract-criterion}
 s|\mathcal P|
 \left(\sum_{j=0}^t\binom kj\right)
 \left(\frac p{s-1}\right)^{(\gamma k-t)/\mu}
 +sT e^{-pk/(8\mu)}<1,
\end{equation}
there is a coloring \(c:X\to\Z/s\Z\) with
\(d_{\rm mono}(c,P)\ge t\) for every \(P\in\mathcal P\).
\end{theorem}

\begin{proof}
For each key take an independent increment \(E_\kappa\in\Z/s\Z\)
with distribution
\[
 \Prb(E_\kappa=0)=1-p,\qquad
 \Prb(E_\kappa=u)=p/(s-1)\quad(u\ne0),
\]
and put \(c(x)=c_0(x)+E_{\kappa(x)}\).
All occurrences of a key use the same increment.

Fix a target \(b\). For a rich tuple let
\(I_b=\{j:c_0(x_j)\ne b\}\), so \(|I_b|\ge\gamma k\).
If the final tuple has fewer than \(t\) positions different from \(b\),
there is a set \(F\subseteq I_b\), \(|F|\le t\), outside which
every position in \(I_b\) changes to \(b\).
This prescribes a particular nonzero increment on at least
\((\gamma k-t)/\mu\) distinct keys.
Incompatible prescriptions on a shared key give probability zero;
compatible prescriptions have probability at most
\((p/(s-1))^{(\gamma k-t)/\mu}\).
There are at most \(\sum_{j=0}^t\binom kj\) choices of \(F\).
Union over tuples and targets gives the first term in
\eqref{eq:abstract-criterion}.

For an exceptional tuple, write
\(S_b=\#\{j:c(x_j)\ne b\}\).
It is a sum of independent contributions \(X_\kappa\in[0,\mu]\).
At a position whose outer color is \(b\), the probability of lying
outside \(b\) is \(p\); at any other position it is
\(1-p/(s-1)\ge p\). Hence
\(\E S_b\ge pk\).

For completeness, if independent \(Y_i\in[0,1]\) have total mean
\(\nu\), convexity gives
\[
 \E e^{-uY_i}\le1-(1-e^{-u})\E Y_i.
\]
Multiplying and using \(1-v\le e^{-v}\) shows that
\(\Prb(\sum Y_i\le\nu/2)\le e^{-\nu/8}\):
take \(u=\log2\), for which the exponent is
\(-\nu(1-\log2)/2\le-\nu/8\).
Apply this to \(Y_\kappa=X_\kappa/\mu\).
The event \(S_b<t\) is contained in
\(\{S_b\le(\E S_b)/2\}\), and so has probability at most
\(e^{-pk/(8\mu)}\).

Two tuples with the same literal signature give exactly the same
event on the same independent increments. They need be counted only
once for each target. The second term of
\eqref{eq:abstract-criterion} therefore bounds all exceptional failures.
If their sum is less than one, a successful choice exists.
\end{proof}

This theorem concerns distance from constant words. For \(s=3\) it does
not guarantee a positive count of every color separately. The distinction
allows the coefficient \(1/18\) below.

\subsection{The exact two-point geometric bound}

On our cyclic group define
\begin{equation}\label{eq:keys}
 f(n)=\norm{q\widetilde y(n)}_2^2,\qquad
 \kappa(n)=\left((V(y_i(n)))_{i=1}^D,\ \fl{f(n)}\right).
\end{equation}
The key remembers a specific box and a specific integer squared-norm
band.

\begin{lemma}[Key multiplicity at most two]\label{lem:multiplicity}
Every key occurs at most twice on any nonzero-step \(k\)-term cyclic
progression.
\end{lemma}

\begin{proof}
Restrict to visits to one second-system box. By
\cref{lem:box-affine}, their centered representatives lie on one real
affine line, even if the visit indices are not consecutive.
Any two distinct samples become separated by Euclidean distance at
least one after multiplication by \(q\).
Indeed, choose a nonzero integer representative \(z\) of
\(\lambda(n_j-n_{j'})\) modulo \(q^D\), and write \(z=q^a u\)
with \(a<D\), \(q\nmid u\).
In coordinate \(a+1\) the circle displacement is \(u/q\) modulo one,
whose distance from zero is at least \(1/q\).
Real representative differences are at least circle distances.
Thus
\begin{equation}\label{eq:separation}
 \norm{q\widetilde y(n_j)-q\widetilde y(n_{j'})}_2\ge1.
\end{equation}
The argument only requires \(q\nmid u\), not that \(u\) be a unit
modulo \(q\).

Parameterize the line at unit speed from its nearest point to the
origin. The squared norm has form \(t^2+c\).
If two sample parameters on the same side satisfy \(0\le t_1<t_2\)
and belong to one half-open band \([m,m+1)\), then
\[
 (t_2-t_1)^2\le t_2^2-t_1^2<1.
\]
This contradicts \eqref{eq:separation}.
There is at most one sample on each side, hence at most two in the
band. A sample at \(t=0\) can only reduce this count.
\end{proof}

The half-open outer boundary makes the last inequality strict.
Replacing it by a closed interval and retaining only its diameter
loses this useful factor of two.
The bound is optimal for arbitrary separated points on a line in a
norm band (take parameters \(-1/2,1/2\) in the band \([0,1)\)).
No claim of attainment by an actual key of our arithmetic embedding
is needed.

\subsection{Affine signatures have two separate entropy costs}

Call a progression affine when all of its centered second-coordinate
representatives have one expression \(A_0+jB_0\) in \(\R^D\).
Its signature is the complete word
\[
 ((c_0(n_j),\kappa(n_j)))_{j<k}.
\]

\begin{lemma}[Affine-signature count]\label{lem:signature}
If \(T_{\rm aff}\) is the number of such signatures, then
\begin{align}
 T_{\rm aff}&\le
 C_3T_{\rm glob}\bigl(k(Dq^2+2)\bigr)^3,
 \label{eq:signature-count}\\
 \log\max\{1,T_{\rm aff}\}
 &\le C_A k^{2/3}L^{5/3}
 \label{eq:signature-entropy}
\end{align}
for a fixed constant \(C_A\) after \(A\) is chosen.
The bound holds uniformly for \(s=2,3\), \(B\), and positive \(C\)
in \eqref{eq:target}, after increasing the onset in \(k\).
\end{lemma}

\begin{proof}
On an affine progression,
\[
 f(n_j)=q^2\norm{A_0+jB_0}_2^2
       =\theta_0+\theta_1j+\theta_2j^2,\qquad
 0\le f(n_j)\le Dq^2/4.
\]
Compare these \(k\) values with every relevant integer boundary.
There are at most \(k(Dq^2+2)\) affine hyperplanes in the three
parameters \(\theta_0,\theta_1,\theta_2\).
Their joint signs determine all integer parts, including exact
integer values. The arrangement bound gives the cubic factor in
\eqref{eq:signature-count}.
The full-label word determines both the specific \(V\)-boxes and all
outer colors, because \(c_*\) has already been fixed.
Multiplying the two upper counts proves \eqref{eq:signature-count};
no independence between them is being asserted.

Taking logarithms gives
\[
 \log\max\{1,T_{\rm aff}\}
 \le \log T_{\rm glob}+6\log q+O(\log k+\log D).
\]
Use \eqref{eq:global-entropy} and \eqref{eq:q-rounding}.
For the stated uniformity, \(B\) is absent from these counts, and the
negative term \(-Ck\) only decreases the upper bound for \(\log q\).
Both \(a_2,a_3\) are bounded absolute constants.
\end{proof}

It is essential to retain \emph{both} \(\log T_{\rm glob}\) and
\(6\log q\). Optimizing only the quadratic part would permit an
unjustified larger exponent.

\subsection{Checking the two failure probabilities}

\begin{theorem}[The robust cyclic endpoint]\label{thm:cyclic}
For common absolute constants \(B,C,K>0\), all integers \(k\ge K\),
and \(s\in\{2,3\}\), there is a coloring of a cyclic group
\(\Z/N_s\Z\) such that
\begin{align}
 \log N_s&\ge
 \frac{s-1}{12s}k(L-5\log L)-Ck,\label{eq:cyclic-size}\\
 d_{\rm mono}(c,P)&\ge
 \fl{\frac B2 k^{2/3}L^{5/3}}
 \quad\text{for every nonzero-step \(k\)-term progression \(P\)}.
 \label{eq:cyclic-robust}
\end{align}
\end{theorem}

\begin{proof}
Choose \(A\) so that \cref{prop:outer-tests} applies, then obtain
\(C_A\) from \cref{lem:signature}.
Fix \(B>32(C_A+1)\), increasing it if necessary to have \(B>2\).
Put \(E_k=k^{2/3}L^{5/3}\); then \(pk=BE_k\).
We will choose \(C\) below and then take \(k\) sufficiently large.

There are fewer than \(N_s^2\) nonzero-step cyclic progressions.
By \cref{thm:dichotomy}, every one is affine or rich with
\[
 \gamma_s=g_s-4/L,\qquad g_s=1-1/s.
\]
By \cref{lem:multiplicity}, key multiplicity is at most two.
Since \(p\to0\) and \(\gamma_s\to g_s>0\), we also have
\(t=\fl{pk/2}<\gamma_s k\) for all sufficiently large \(k\).
Apply \cref{thm:abstract} with \(t=\fl{pk/2}\) and \(T=T_{\rm aff}\).
The affine term is at most
\[
 sT_{\rm aff}e^{-pk/16}
 \le \exp\!\left[\log3+(C_A-B/16)E_k\right]=o(1).
\]
In particular \(B\) is a fixed constant chosen to dominate entropy;
there is no need for that entropy to be \(o(pk)\).

For the rich term, let
\(\mathcal B_k(t)=\sum_{j=0}^t\binom kj\).
For \(1\le t<k\), the elementary estimate
\[
 \mathcal B_k(t)\le (t+1)(ek/t)^t
\]
holds because \(\binom kj\le(ek/j)^j\) and the latter expression
increases with \(j\le k\).
It follows that
\begin{equation}\label{eq:exception-entropy}
 \log\mathcal B_k(t)=O(k^{2/3}L^{8/3})=o(k),
 \qquad
 t\,\abs{\log(p/(s-1))}=o(k).
\end{equation}
Thus the logarithm of the rich term is at most
\begin{equation}\label{eq:rich-log}
 \log s+2\log N_s+\log\mathcal B_k(t)
 +\frac{\gamma_s k-t}{2}\log\frac p{s-1}.
\end{equation}
By the prime-power rounding bound,
\(\log N_s=Q_s(k)+o(k)\), and
\[
 \log\frac p{s-1}
 =-\frac L3+\frac53\log L+\log\frac B{s-1}.
\]
Since \(a_s=g_s/12\), the terms of size \(kL\) and \(k\log L\)
in \eqref{eq:rich-log} cancel exactly:
\begin{align}
 2a_s-\frac{g_s}{6}&=0,&
 -10a_s+\frac{5g_s}{6}&=0.\label{eq:cancellations}
\end{align}
The remaining upper bound is
\begin{equation}\label{eq:rich-linear}
 \left[-2C+\frac{g_s}{2}\log\frac B{s-1}
                 +\frac23+o(1)\right]k.
\end{equation}
The term \(2/3\) comes from the loss \(4/L\) in the rich proportion.
The exception terms and rounding are \(o(k)\) by
\eqref{eq:exception-entropy}.
For example, choosing \(C=2+\frac14\log B\) makes the coefficient
negative for both \(s=2,3\). The rich term is consequently \(o(1)\).

For all sufficiently large \(k\), the two failure bounds sum to less
than one. \Cref{thm:abstract} supplies the coloring.
Finally \(D\log q\ge Q_s(k)\), giving \eqref{eq:cyclic-size}.
Every choice of constant preceded the final threshold, and only the
two palettes were involved, so a common absolute \(K\) suffices.
\end{proof}

\subsection{A finite description of the colorings}

\begin{corollary}[Subpower description size]\label{cor:description}
The cyclic colorings in \cref{thm:cyclic} have descriptions of length
\[
 \exp\!\left(O(k^{2/3}(\log k)^{5/3})\right)=N_s^{o(1)}
\]
bits, allowing the outer-color table, the increment table, and the
integer parameters of the construction.
\end{corollary}

\begin{proof}
There are
\(|\cL|=|U|^D|V|^D\) outer labels, with
\(\log|\cL|=O_A(E_k)\).
There are at most \(|V|^D(\fl{Dq^2/4}+1)\) possible keys, also with
logarithm \(O_A(E_k)\).
Store their colors and increments in a fixed lexicographic ordering,
using a fixed dummy value on unused entries.
Even storing an explicit index for each used entry changes the
exponential bound only by a polynomial factor in the logarithmic
parameters.
The finite mesh is specified by its integer \(J_0\) and rational
\(\alpha\): indeed
\[
 a_j=\alpha\left[
 \left(1+\frac1{2\alpha}\right)^{j/J_0}-1\right].
\]
Its endpoints therefore have finite algebraic descriptions.
Since \(\log N_s=\Theta(kL)\) and \(E_k=o(kL)\), the claimed
\(N_s^{o(1)}\) relation follows.
\end{proof}

This is a bound on description length. It does not provide an efficient
algorithm for finding successful tables or for certifying all local
tests.
